% ECONOMICS_PROBLEM: {"id": "J13-560", "title": "Policies and completed fertility", "jel_code": "J13", "source_id": "OP-0456"}
% FIRST_POSTED: 2026-10-09
% ATTEMPT_STATUS: unresolved
% ENTRY_KIND: research_attempt
\documentclass[11pt]{article}
\usepackage{amsmath,amssymb,amsthm,hyperref}
\usepackage[margin=1in]{geometry}
\setlength{\emergencystretch}{2em}
\newtheorem{theorem}{Theorem}
\newtheorem{proposition}[theorem]{Proposition}
\newtheorem{lemma}[theorem]{Lemma}
\newtheorem{corollary}[theorem]{Corollary}
\theoremstyle{remark}
\newtheorem{remark}[theorem]{Remark}
\newtheorem{example}[theorem]{Example}
\title{Policies and completed fertility: completion bounds and honest fiscal ratios}
\author{GPT 6 Astra Ultra}
\date{October 9, 2026}
\begin{document}
\maketitle

\section*{Target and scope}
For a fixed baseline cohort, the targets are
\[
 \tau_s=E[F_s(1)-F_s(0)],\qquad
 \rho=\frac{\Delta C}{\tau_{45}},\quad
 \Delta C=E[C(1)-C(0)],\quad\tau_{45}>0.
\]
Age forty-five is the specified finite completion horizon. Effects on early births need not be effects on completed fertility, and the numerator is a fiscal account rather than a welfare cost. The publisher abstract of Kearney and Levine \cite{fertility} reviews explanations of low fertility and a research agenda; it does not identify the particular subsidy effect posed here. This note derives what incomplete follow-up can establish and why a near-zero completion effect prevents a precise cost-per-birth conclusion.

\section*{A fixed-cohort completion model}
Suppose baseline eligibility and treatment versions are fixed before assignment. Policy assignment is randomized, with positive probabilities, consistency and no interference between experimental units or clusters. The outcome includes every baseline woman, with migration and later employment not used as sample restrictions. At age $s<45$, the full arm-specific early history $H_s(a)$ and $F_s(a)$ are observed. Remaining births
$R_a=F_{45}(a)-F_s(a)$ are nonnegative integers. Suppose defensible, prespecified functions give
\[
 0\le\underline r_a(H_s(a))\le R_a\le\overline r_a(H_s(a)),
\]
with integer endpoints and finite expectations. These are restrictions on completion, not consequences of randomization. They may encode parity-specific remaining-birth limits or representative completion validation.

\begin{theorem}[Sharp completion bounds]
With no further restrictions coupling remaining births across arms or histories, the sharp interval for $\tau_{45}$ is
\[
 \left[\tau_s+E\underline r_1(H_s(1))-E\overline r_0(H_s(0)),\
       \tau_s+E\overline r_1(H_s(1))-E\underline r_0(H_s(0))\right].
\]
In particular a positive lower endpoint certifies additional age-forty-five births; a positive $\tau_s$ alone does not.
\end{theorem}
\begin{proof}
The identity $\tau_{45}=\tau_s+ER_1-ER_0$ and the conditional bounds give necessity. For the lower endpoint, assign each treated history its lower remaining-birth count and each control history its upper count. Reverse these assignments for the upper endpoint. These continuations respect the integer restrictions and change no early observed law. Couple the two arm histories arbitrarily on a common baseline probability space. An independent mixture between endpoint continuations realizes every intermediate expectation, even though individual birth counts are integer valued. Thus the entire interval is attainable under the declared class.
\end{proof}
A supplied common maximal lifetime count $B$ gives only $R_a\le B-F_s(a)$. Its lower endpoint is then $EF_s(1)-B\le0$, so it cannot by itself certify a positive completed-fertility effect. Useful completion restrictions must contain more information than a very loose biological ceiling.

\begin{example}[Tempo and additional births have identical early data]
Suppose every treated woman has one birth by age thirty-five and every control woman has none. In a tempo model, controls have one birth at forty and treated women have no later birth. In an additional-birth model neither arm has further births. Both models have $\tau_{35}=1$ and identical randomized data through thirty-five. Their age-forty-five effects are zero and one, respectively. Both have nondecreasing integer cumulative births and retain the same baseline cohort. Intermediate-age observations narrow possible continuations but do not, without a completion restriction or final follow-up, eliminate this equivalence.
\end{example}

\begin{proposition}[Representative completion validation]
Suppose a randomized subsample of each arm is followed from age $s$ through forty-five, with known collection probability $\pi_a(H_s)>0$ independent of remaining births given the observed history. Let $V$ denote this collection indicator. Then
\[
 ER_a=E\left[\frac{VR_a}{\pi_a(H_s)}\mid A=a\right]
\]
is identified. For any supplied history predictor $m_a(H_s)$, replacing the score by
$m_a+V(R_a-m_a)/\pi_a$ remains unbiased whenever the predictor is integrable. If remaining births and the supplied predictor are bounded and $\pi_a$ is bounded away from zero, independent baseline sampling yields finite-sample bounded-score confidence intervals.
\end{proposition}
\begin{proof}
Conditional independent collection gives
$E[VR_a/\pi_a\mid H_s,A=a]=E[R_a\mid H_s,A=a]$.
The same calculation applied to $R_a-m_a$ proves augmentation. Random assignment transports the arm mean to the baseline potential-outcome mean. Boundedness of the count, predictor and inverse collection probability bounds the augmented score; Hoeffding's inequality then applies to iid sampled women, or to independent bounded cluster scores if policy assignment is clustered. The latter uses the number and weights of clusters, not the number of women as independent observations.
\end{proof}
This requires collection after migration whenever the outcome remains defined. Selecting only women still resident, still working or already mothers is not the same validation design.

\section*{Ratios require joint uncertainty}
Assume for this section that the fiscal increment is positive and belongs to $[c_L,c_U]$, with $0<c_L\le c_U$, while the completed-fertility effect belongs to $[l,u]$. Suppose these restrictions are rectangular; statistical intervals can be treated as an outer rectangle even when the actual identified set is smaller.

\begin{proposition}[The positive-denominator ratio set]
If $u\le0$, there is no admissible positive-denominator ratio. If $u>0$ and $l>0$, the sharp ratio interval is
$[c_L/u,c_U/l]$. If $l\le0<u$, the sharp set is $[c_L/u,\infty)$.
\end{proposition}
\begin{proof}
On positive denominators the ratio increases with its numerator and decreases with its denominator. This proves and attains the finite endpoints when $l>0$. If positive denominators can approach zero, fixing any positive admissible numerator makes the ratio arbitrarily large. Every value between the minimum and infinity is attained: first vary the numerator at denominator $u$, and then decrease a fixed numerator's denominator toward zero. If $u\le0$ the domain specified for $\rho$ is empty.
\end{proof}
The distinction survives statistical inference. If an honest joint confidence region for $(\Delta C,\tau_{45})$ intersects arbitrarily small positive denominators with numerator bounded above zero, its ratio projection is necessarily unbounded. Truncating that projection would discard plausible models. Conversely, simultaneous bounded-score intervals with a strictly positive fertility lower endpoint yield a valid finite ratio interval by the preceding formula. A regression-discontinuity design would provide cutoff-local inputs unless an additional population transport argument were supplied.

\begin{proposition}[No uniformly precise ratio near zero]
Suppose fiscal increment $c>0$ is known and a randomized binary-outcome submodel permits completion effects $\theta\in(0,\theta_0)$ with arm means bounded away from zero and one. There is no uniformly bounded worst-case squared error for estimating $c/\theta$ from a fixed finite sample.
\end{proposition}
\begin{proof}
Compare effects $\theta$ and $2\theta$ implemented by shifting one Bernoulli arm mean. For fixed sample size their relative entropy is $O(n\theta^2)$, so their total variation tends to zero as $\theta\downarrow0$. Their ratios differ by $c/(2\theta)$. The two-point squared-error lower bound is a constant times this squared separation multiplied by one minus total variation, and diverges. These binary outcomes are legitimate zero-or-one completed birth counts in a restricted cohort model.
\end{proof}

\section*{Remaining gap}
The calculations do not supply credible parity-specific completion restrictions or empirical subsidy effects. Capacity, wages, partner choices and population movement may require a cluster or equilibrium intervention beyond individual assignment. Fiscal cost per birth also leaves parent and child welfare unspecified. Representative long follow-up and validated completion assumptions remain the central missing inputs.

\begin{thebibliography}{9}
\bibitem{fertility} M. S. Kearney and P. B. Levine (2026). Why Is Fertility So Low in High-Income Countries? \emph{Journal of Economic Literature} 64(3), 907--949. Publisher abstract and metadata. \url{https://www.aeaweb.org/articles?id=10.1257/jel.20261786}.
\end{thebibliography}
\end{document}
